\documentclass[a4print,english,lof,lot,twoside]{univpmthesis}
\errorcontextlines=9

\RequirePackage[utf8]{inputenc}
\RequirePackage[T1]{fontenc}

\usepackage{lmodern}

%%%%%%%%%%%%%%%%%%%%%%%%%%%%%%%%%%%%%%%%%%%%%%%%%%%%%%%%%%%
% Metadata
%%%%%%%%%%%%%%%%%%%%%%%%%%%%%%%%%%%%%%%%%%%%%%%%%%%%%%%%%%%
\thfaculty{Facolt\`{a} di Ingegneria}
\thcourse{Corso di Laurea in Ingegneria Meccanica}
\thtitle{Titolo della Tesi}
\thsubtitle{Sottotitolo della Tesi} % NON NECESSARIO
\thauthor{Nome Studente}
\thadvisor{Prof.~ Nome Relatore}
\thcoadvisor{Prof.~Nome Correlatore 1}{Prof.~Nome Correlatore 2}{}


\ayear{20XX-20XX}
\thesisdedication{A tutti gli Studenti dell'UNIVPM\\ per i loro primi passi in \LaTeX}
\thlocation{Ancona}
\thtime{Gennaio 20XX} % data scrittura
%%%%%%%%%%%%%%%%%%%%%%%%%%%%%%%%%%%%%%%%%%%%%%%%%%%%%%%%%%%

\usepackage{amsmath,amsfonts,amssymb,bm} %<- scrittura equazioni
%\usepackage[italian,english]{babel}
\usepackage{graphicx} 	%<- immagini
\usepackage{epstopdf} 	%<- conversione immagini .eps
\usepackage{subfig} 	%<- sottofigure
\usepackage{footnote}	%<- note a piè pagina
\usepackage{tabularx,booktabs,multicol,multirow} %<- tabelle
\usepackage{mathrsfs}
\usepackage{caption}	%<- didascalie
\usepackage{microtype}
\usepackage{gensymb,siunitx}	%<- simboli ed unità di misura
\usepackage{hyperref}
\usepackage{glossaries}
\usepackage{verbatim}
\usepackage{xcolor}
\usepackage[a-1b]{pdfx} %<- generazione pdf-A

% Solo per generare testo, da cancellare quando si inizia a scrivere:
\usepackage{lipsum}

%%%%%%%%%%%%%%%%%%%%%%%%%%%%%%%%%%%%%%%%%%%%%%%%%%%%%%%%%%%
\newcommand{\eqname}{Eq.}
%%%%%%%%%%%%%%%%%%%%%%%%%%%%%%%%%%%%%%%%%%%%%%%%%%%%%%%%%%%

\begin{document}

%%%%%%%%%%%%%%%%%%%%%%%%%%%%%%%%%%%%%%%%%%%%%%%%%%%%%%%%%%%
% Front matter contents
%%%%%%%%%%%%%%%%%%%%%%%%%%%%%%%%%%%%%%%%%%%%%%%%%%%%%%%%%%%
\frontmatter

\maketitle

\begin{thesisacknowledge}[italian]
	 Una menzione particolare va a Marco Giammarini, primo sviluppatore della classe \texttt{univpmphdthesis} per la scrittura della Tesi di Dottorato per la Facoltà di Ingegneria dell'Università Politecnica delle Marche, e da cui la presente classe per Lauree Triennali e Magistrali si ispira.
	 
	 Si ringrazia per il prezioso supporto il forum \emph{GuIT}, \emph{TeX Stack Exchange} e \emph{StackOverflow}, luci nel buio dei nostri errori di compilazione.
\end{thesisacknowledge}

\begin{thesisabstract}[english]
	The \texttt{univpmthesis} class produces a template for the Bachelor and Master thesis manuscripts for students at UNIVPM. The class allows to write your manuscript using either italian or english; so, if you are an international student, feel free to use it!
\end{thesisabstract}

\begin{thesisabstract}[italian]
	La classe \texttt{univpmthesis} fornisce un template per la scrittura della tesi triennale e Magistrale per studenti dell'UNIVPM. La classe permette di impiegare sia la lingua inglese che l'italiano, sentiti libero di utilizzarla come preferisci!
\end{thesisabstract}

\thesistoc

%%%%%%%%%%%%%%%%%%%%%%%%%%%%%%%%%%%%%%%%%%%%%%%%%%%%%%%%%%%
% Main matter contents
%%%%%%%%%%%%%%%%%%%%%%%%%%%%%%%%%%%%%%%%%%%%%%%%%%%%%%%%%%%
\mainmatter

\input{chapters/chapter1/chapter1}

%%%%%%%%%%%%%%%%%%%%%%%%%%%%%%%%%%%%%%%%%%%%%%%%%%%%%%%%%%%
% Back matter contents
%%%%%%%%%%%%%%%%%%%%%%%%%%%%%%%%%%%%%%%%%%%%%%%%%%%%%%%%%%%
\backmatter

%\bibliographystyle{IEEEbib}
%\bibliography{IEEEabrv,bibliografia}

\bibliographystyle{unsrt}
\bibliography{frontbackmatters/bibliography.bib}

\end{document}
